\documentclass[11pt]{article}
\usepackage{mathblatt}
\usepackage{adjustbox,varwidth}
% gesetzt von werkzeuge/setzer.py aus dem Lernweg (L5-A · L5-B · L5-C · L5-D · L5-T) und den Bankzeilen; Muster T6B.
% Präambel des Setzers (werkzeuge/setzer.py), wortgleich aus dem Hand-Satz T6B (bau/proben/2026-10-09/kathete-berechnen), dort aus K4W.
\makeatletter
\setlength{\mbpfbe}{0mm}% keine Punkte (Bauregel 6.4)
% eingerückt (Bauregel 4.7, 6.6): \gEin Einzug, \gSchrift eine Stufe kleiner
\newlength{\gEin}\setlength{\gEin}{0pt}
\let\gSchrift\relax
\renewcommand{\pfkreuzzeile}[1]{\par\noindent\hangindent3.2em\hangafter1\hspace*{1.6em}$\square$\ #1\par}
\newcommand{\gkreuz}[1]{$\square$\ #1\hspace{2em}}
% Bauregel 6.7: links, was man ansieht, rechts, was man tut; Spaltengrenze überall an derselben Stelle
\newsavebox{\gbox}\newlength{\gLinks}
\newcommand{\gzwei}[2]{\par\smallskip\noindent\sbox{\gbox}{#1}%
  \setlength{\gLinks}{\dimexpr0.5\textwidth-\gEin-\mbpfnr\relax}%
  \ifdim\wd\gbox>\dimexpr\gLinks-4mm\relax
    \usebox{\gbox}\par\smallskip\noindent\begin{minipage}{\linewidth}#2\end{minipage}%
  \else
    \begin{minipage}[t]{\gLinks}\vspace{0pt}\usebox{\gbox}\end{minipage}%
    \begin{minipage}[t]{\dimexpr\linewidth-\gLinks\relax}\vspace{0pt}\raggedright #2\end{minipage}%
  \fi\par}
\RenewDocumentEnvironment{pfaufg}{m m m}{%
  \begin{lrbox}{\mbaufgabebox}\begin{minipage}[t]{\linewidth}%
  \gSchrift\noindent\hspace*{\gEin}\mbpfrandmarke{#1}%
  \begin{minipage}[t]{\mbpfnr}\strut\textbf{#2}\end{minipage}%
  \begin{minipage}[t]{\dimexpr\linewidth-\gEin-\mbpfnr\relax}\raggedright\strut\ignorespaces}%
 {\end{minipage}%
  \end{minipage}\end{lrbox}%
  \setlength{\mbaufgabehoehe}{\dimexpr\ht\mbaufgabebox+\dp\mbaufgabebox\relax}%
  \par\addvspace{7pt}\Needspace*{\mbaufgabehoehe}%
  \noindent\usebox{\mbaufgabebox}\par\addvspace{7pt}}
\makeatother
% Rechenraum als Karo (Bauregel 6.9): n Rechenschritte = n mal 10 mm
\newcommand{\karo}[1]{\par\smallskip\noindent\begin{tikzpicture}
  \clip (0,0) rectangle ({\linewidth-1mm},{#1*10mm});
  \draw[black!16,line width=0.3pt,step=5mm] (0,0) grid ({\linewidth},{#1*10mm});
  \end{tikzpicture}\par}
\newcommand{\kk}[1]{$\square$\,#1\hspace{0.9em}}
\newcommand{\linie}[1]{\,{\color{black!45}\rule[-1pt]{#1}{0.4pt}}\,}
% Zeile am Ende jedes Durchgangs: Originale klein grau, darüber; dann Vorgänger/Nachfolger (Bauregel 3.4, 6.5)
\newcommand{\blattende}[2]{\par\vfill\noindent\begin{minipage}{\linewidth}{\scriptsize\color{mbgrau}Originale: #1\par}\smallskip
{\footnotesize #2\par}\end{minipage}}
\newcommand{\pkt}{\textperiodcentered{}}

\newcommand{\kennung}{MZE}
\newcommand{\grau}[1]{{\color{mbgrau}#1}}

\begin{document}
\pfheftstil
\fancyfoot[C]{\parbox[t]{\textwidth}{\mbfhbox\par\vspace{1.5mm}{\footnotesize\color{mbgrau}{\scriptsize \kennung}\hfill\thepage}}}
\pfheftkopf{Termwerte berechnen}{Klasse 7}
% L5-A: Zahl einsetzen und ausrechnen
\begin{pfaufg}{}{1.}{}
\grau{a)\ Berechne den Wert des Terms $3x + 2 \cdot (x - 1)$ für $x = 4$.\par\smallskip Für jedes $x$ die $4$ einsetzen; aus $3x$ wird $3 \cdot 4$: \quad $3 \cdot 4 + 2 \cdot (4 - 1)$\par Klammer zuerst: \quad $4 - 1 = 3$, \ also \ $3 \cdot 4 + 2 \cdot 3$\par Punkt vor Strich: \quad $12 + 6 = 18$\par Der Termwert ist $18$.}
\par\medskip
b)\ Berechne den Wert von $5x - 4$ für \ a)~$x = 2$ \qquad b)~$x = 3$ \qquad c)~$x = 10$.\karo{3}
\fusshilfe{\mbox{1:~b) $6;\ 11;\ 46$}}
\end{pfaufg}

% L5-A: Zahl einsetzen und ausrechnen
\begin{pfaufg}{}{2.}{}
Berechne den Wert für $x = 3$. \par\vspace{3pt} a)~$7 + 2x$ \qquad b)~$4 \cdot (x + 2)$ \qquad c)~$20 - 3x$ \qquad d)~$2x + x - 1$\karo{3}
\fusshilfe{\mbox{2:~$13;\ 20;\ 11;\ 8$}}
\end{pfaufg}

% L5-A: Zahl einsetzen und ausrechnen
\begin{pfaufg}{}{3.}{}
\gzwei{\adjustbox{max width=\linewidth}{\begin{varwidth}{50cm}\wertetabelle{x}{3x - 2}{1,2,3,4,5}\end{varwidth}}}{%
Berechne die Termwerte von $3x - 2$ und trage sie in die Tabelle ein.\karo{3}}
\fusshilfe{\mbox{3:~$1;\ 4;\ 7;\ 10;\ 13$}}
\end{pfaufg}

% L5-A: Zahl einsetzen und ausrechnen
\begin{pfaufg}{}{4.}{}
Ein Fahrradverleih berechnet den Preis in Euro mit dem Term $4 + 3h$ ($h$: Anzahl der Stunden). Tom hat 20\,€ und möchte das Rad 6 Stunden leihen. Reicht sein Geld? Wie viel bleibt übrig oder fehlt?\karo{3}
\fusshilfe{\mbox{4:~nein, $22$\,€; es fehlen $2$\,€}}
\end{pfaufg}

% L5-B: Negative Zahlen einsetzen
\begin{pfaufg}{}{5.}{}
\grau{a)\ Berechne den Wert des Terms $7 - 3 \cdot (x + 2)$ für $x = -5$.\par\smallskip $-5$ einsetzen: \quad $7 - 3 \cdot (-5 + 2)$\par Klammer zuerst: \quad $-5 + 2 = -3$, \ also \ $7 - 3 \cdot (-3)$\par Punkt vor Strich: \quad $3 \cdot (-3) = -9$, \ also \ $7 - (-9)$\par Minus vor Minus wird Plus: \quad $7 + 9 = 16$}
\par\medskip
b)\ Berechne den Wert von $3x + 10$ für \ a)~$x = -2$ \qquad b)~$x = -5$ \qquad c)~$x = -1$.\karo{3}
\fusshilfe{\mbox{5:~b) $4;\ -5;\ 7$}}
\end{pfaufg}

% L5-B: Negative Zahlen einsetzen
\begin{pfaufg}{}{6.}{}
Berechne den Wert für $x = -3$. \par\vspace{3pt} a)~$8 - x$ \qquad b)~$-2x$ \qquad c)~$5 - 4x$ \qquad d)~$x - 6$\karo{3}
\fusshilfe{\mbox{6:~$11;\ 6;\ 17;\ -9$}}
\end{pfaufg}

% L5-B: Negative Zahlen einsetzen
\begin{pfaufg}{}{7.}{}
Berechne. \par\vspace{3pt} a)~$6 \cdot (x - 1) + 20$ für $x = -2$ \qquad b)~$10 - 2 \cdot (x + 3)$ für $x = -7$\karo{3}
\fusshilfe{\mbox{7:~$2;\ 18$}}
\end{pfaufg}

% L5-B: Negative Zahlen einsetzen
\begin{pfaufg}{}{8.}{}
Welcher Term hat für $x = -2$ den größeren Wert: $3 - 4x$ oder $2 \cdot (x + 6)$? Um wie viel ist er größer?\karo{3}
\fusshilfe{\mbox{8:~$3 - 4x$ (Wert $11$), um $3$ größer}}
\end{pfaufg}

% L5-C: Quadrate
\begin{pfaufg}{}{9.}{}
\grau{a)\ Berechne den Wert des Terms $2x^2 - 3x$ für $x = -4$.\par\smallskip $-4$ in Klammern einsetzen: \quad $2 \cdot (-4)^2 - 3 \cdot (-4)$\par Hochzahl zuerst: \quad $(-4)^2 = (-4) \cdot (-4) = 16$, \ also \ $2 \cdot 16 - 3 \cdot (-4)$\par Punkt vor Strich: \quad $2 \cdot 16 = 32$ \ und \ $3 \cdot (-4) = -12$, \ also \ $32 - (-12)$\par Minus vor Minus wird Plus: \quad $32 + 12 = 44$}
\par\medskip
b)\ Berechne den Wert von $x^2$ für \ a)~$x = 5$ \qquad b)~$x = -5$ \qquad c)~$x = -1$ \qquad d)~$x = -10$.\karo{3}
\fusshilfe{\mbox{9:~b) $25;\ 25;\ 1;\ 100$}}
\end{pfaufg}

% L5-C: Quadrate
\begin{pfaufg}{}{10.}{}
Berechne den Wert für $x = -2$. \par\vspace{3pt} a)~$4x^2$ \qquad b)~$x^2 + 4$ \qquad c)~$x^2 - x$\karo{3}
\fusshilfe{\mbox{10:~$16;\ 8;\ 6$}}
\end{pfaufg}

% L5-C: Quadrate
\begin{pfaufg}{}{11.}{}
Berechne. \par\vspace{3pt} a)~$10 - x^2$ für $x = -3$ \qquad b)~$2x^2 + 5x$ für $x = -1$\karo{3}
\fusshilfe{\mbox{11:~$1;\ -3$}}
\end{pfaufg}

% L5-C: Quadrate
\begin{pfaufg}{}{12.}{}
\gzwei{\adjustbox{max width=\linewidth}{\begin{varwidth}{50cm}\wertetabelle{x}{x^2 - 2x}{-2,-1,0,1,2,3}\end{varwidth}}}{%
Berechne die Termwerte von $x^2 - 2x$ und trage sie in die Tabelle ein. Für welche Zahl $x$ ist der Termwert negativ?\karo{3}}
\fusshilfe{\mbox{12:~$8;\ 3;\ 0;\ -1;\ 0;\ 3$}}
\end{pfaufg}

% L5-D: Zwei Variablen und Bruchstrich
\begin{pfaufg}{}{13.}{}
\grau{a)\ Berechne den Wert des Terms $\dfrac{a + b}{c}$ für $a = 3$, $b = -9$ und $c = -2$.\par\smallskip Jede Variable durch ihre Zahl ersetzen: \quad $\dfrac{3 + (-9)}{-2}$\par Oben und unten zuerst: \quad Zähler $3 + (-9) = -6$, \ Nenner $-2$\par Teilen: \quad $(-6) : (-2) = 3$}
\par\medskip
b)\ Berechne für $a = 4$ und $b = -3$. \par\vspace{3pt} a)~$a + b$ \qquad b)~$2a - b$ \qquad c)~$a \cdot b$ \qquad d)~$b - a$\karo{3}
\fusshilfe{\mbox{13:~b) $1;\ 11;\ -12;\ -7$}}
\end{pfaufg}

% L5-D: Zwei Variablen und Bruchstrich
\begin{pfaufg}{}{14.}{}
Berechne den Wert des Terms $\dfrac{a - b}{c}$ für $a = 1$, $b = -7$ und $c = -4$.\karo{3}
\fusshilfe{\mbox{14:~$-2$}}
\end{pfaufg}

% L5-D: Zwei Variablen und Bruchstrich
\begin{pfaufg}{}{15.}{}
Berechne den Wert des Terms $\dfrac{a + b}{c}$ für $a = 6$, $b = -1$ und $c = -2$.\karo{3}
\fusshilfe{\mbox{15:~$-2{,}5$}}
\end{pfaufg}

% L5-D: Zwei Variablen und Bruchstrich
\begin{pfaufg}{}{16.}{}
Für $a = -2$ und $b = 4$: Welcher Term hat den größten Wert? Kreuze an.\par\smallskip \gkreuz{$a \cdot b$}\gkreuz{$\dfrac{b}{a}$}\par \gkreuz{$b - a$}\gkreuz{$a - b$}
\fusshilfe{\mbox{16:~$b - a$ (Wert $6$)}}
\end{pfaufg}

% L5-T: Probetest
\begin{pfaufg}{}{17.}{}
Berechne den Wert von $6 - 2x$ für $x = -4$.\karo{3}
\fusshilfe{\mbox{17:~$14$}}
\end{pfaufg}

% L5-T: Probetest
\begin{pfaufg}{}{18.}{}
Welchen Wert hat der Term $4 \cdot (x - 2)$ für $x = -3$? Kreuze an.\par\smallskip \gkreuz{$-20$}\gkreuz{$-14$}\par \gkreuz{$-4$}\gkreuz{$20$}
\fusshilfe{\mbox{18:~$-20$}}
\end{pfaufg}

% L5-T: Probetest
\begin{pfaufg}{}{19.}{}
Berechne den Wert von $2x^2 + x$ für $x = -3$.\karo{3}
\fusshilfe{\mbox{19:~$15$}}
\end{pfaufg}

% L5-T: Probetest
\begin{pfaufg}{}{20.}{}
Berechne den Wert des Terms $\dfrac{a - b}{c}$ für $a = -1$, $b = 5$ und $c = -4$.\karo{3}
\fusshilfe{\mbox{20:~$1{,}5$}}
\end{pfaufg}

% L5-T: Probetest
\begin{pfaufg}{}{21.}{}
Berechne den Wert von $a^2 - b$ für $a = -3$ und $b = -5$.\karo{3}
\fusshilfe{\mbox{21:~$14$}}
\end{pfaufg}

% L5-T: Probetest
\begin{pfaufg}{}{22.}{}
Fahrschulen rechnen mit dieser Faustformel für den Anhalteweg in Metern ($v$: Geschwindigkeit in km/h): \par\vspace{3pt}\hspace*{15mm}$3 \cdot \dfrac{v}{10} + \left(\dfrac{v}{10}\right)^2$\par\vspace{3pt} 20\,m vor einem Auto läuft ein Kind auf die Straße. Kommt das Auto vor dem Kind zum Stehen, wenn es a)~30\,km/h fährt, b)~50\,km/h fährt?\karo{3}
\fusshilfe{\mbox{22:~a) $18$\,m, ja \quad b) $40$\,m, nein}}
\end{pfaufg}

\par\vfill\noindent{\footnotesize Vorher: Rechnen mit negativen Zahlen\quad\textbullet\quad Weiter: Terme aufstellen\par}
\end{document}
